(1) 各向同性无质量光子满足压强为能量密度的三分之一，故 \(\boxed{p=U/(3V)=aT^4/3}\)。光子数不是守恒热力学变量、化学势为零，由Euler关系 \(TS=U+pV\) 得 \(\boxed{S=4aVT^3/3}\)。

(2) 对 \(U=aVT^4\) 微分，\(dU=aT^4dV+4aVT^3dT\)。由 \(TdS-pdV\) 右侧，\(T d(4aVT^3/3)-aT^4dV/3=aT^4dV+4aVT^3dT\)，与之精确一致。自由能 \(F=U-TS=-aVT^4/3=-pV\)。

(3) 绝热可逆 \(dS=0\) 意味着 \(VT^3=\mathrm{const}\)，故 \(\boxed{TV^{1/3}=\mathrm{const}}\)。总能 \(U=aVT^4\propto V^{-1/3}\)，压强 \(p\propto V^{-4/3}\)。膨胀时光子频率受边界移动红移，辐射温度随体积增加而下降。

(4) 绝热过程中 \(W=\int_{V_i}^{V_f}p,dV=U_i-U_f=aT_i^4V_i[1-(V_i/V_f)^{1/3}]\)。因此辐射气体的绝热指数为 \(4/3\)，而经典单原子理想气体为 \(5/3\)。黑体壁允许吸收与发射光子，固定光子数并非此热平衡模型的约束。
